\documentclass{article}
\usepackage{../mathnotezh}

\title{共点直线系方程的几何意义}
\author{tbj}
\date{}

\begin{document}

\maketitle

设直线$l_1: A_1x + B_1y + C_1 = 0, l_2: A_2x + B_2y + C_2 = 0$.
不妨设$A_1^2 + B_1^2 = A_2^2 + B_2^2 \neq 0$.
过两直线交点$O$的直线的方程为$l: \lambda (A_1x + B_1y + C_1) + \mu (A_2x + B_2y + C_2) = 0$,
这个方程称为过两直线$l_1, l_2$交点的共点直线系方程.
考虑$l$不与$l_1, l_2$重合的情形,此时$\lambda, \mu \neq 0$.

设$l$与$l_1, l_2$的夹角为$\alpha, \beta$,设$l$上一点$P$到$l_1, l_2$的距离分别为$d_1, d_2$,
则有$d_1 = OP \sin \alpha, d_2 = OP \sin \beta$.因此$\dfrac{d_1}{d_2} = \dfrac{\sin \alpha}{\sin \beta}$为定值.

设$P(x_0, y_0)$,则

\[ \lambda (A_1x_0 + B_1y_0 + C_1) + \mu (A_2x_0 + B_2y_0 + C_2) = 0 \]

因此

\[ \lrvert{\lambda} \lrvert{A_1x_0 + B_1y_0 + C_1} =  \lrvert{\mu} \lrvert{A_2x_0 + B_2y_0 + C_2} \]

又因为

\[A_1^2 + B_1^2 = A_2^2 + B_2^2\]

由点到直线的距离公式知

\[\lrvert{\lambda} d_1 = \lrvert{\mu} d_2\]

进而得到

\[\lrvert{\lambda} \sin \alpha = \lrvert{\mu} \sin \beta\]

这样,共点直线系方程的几何意义就是:$l$分$l_1, l_2$的夹角的正弦的比为$\lrvert{\dfrac{\mu}{\lambda}}$.

\begin{rqq}
    注意到,满足$\lrvert{\lambda} \sin \alpha = \lrvert{\mu} \sin \beta$的直线有两条,分别穿过$l_1$与$l_2$所夹的锐角与钝角.
\end{rqq}

\end{document}